\documentclass[aspectratio=1610]{beamer}
\usepackage[utf8]{inputenc}
\usepackage[T1]{fontenc}
\usepackage{xcolor}

\title{Probiers mal mit Gemütlichkeit!}
\subtitle{Vorstellung der Masterthesis}
\date{\today}
\author{Balu der Bär}
\institute{Universität Klagenfurt}
\date{\today}

\usetheme[sidebar=right]{aau}

\begin{document}
	
\begin{frame}
\titlepage
\end{frame}

\begin{frame}
	\tableofcontents
\end{frame}

\section{Freundschaften}
\begin{frame}
	Hallo Freunde!
\end{frame}

\section{Aufzählungen}
\begin{frame}{Ungeordnete Aufzählungen}
	\begin{itemize}
		\item Erster Punkt
		\item Zweiter Punkt
		\begin{itemize}
			\item Erster Unterpunkt
			\item Zweiter Unterpunkt
			\begin{itemize}
				\item Erster Unterunterpunkt
				\item Zweiter Unterunterpunkt
			\end{itemize}
		\end{itemize}
	\end{itemize}
\end{frame}

\begin{frame}{Geordnete Aufzählungen}
	\begin{enumerate}
		\item Erster Punkt
		\item Zweiter Punkt
		\begin{enumerate}
			\item Erster Unterpunkt
			\item Zweiter Unterpunkt
			\begin{enumerate}
				\item Erster Unterunterpunkt
				\item Zweiter Unterunterpunkt
			\end{enumerate}
		\end{enumerate}
	\end{enumerate}
\end{frame}

\section{Alert, Examples, etc.}
\begin{frame}
	Hier kommt ein \alert{alert} und hier kommt eine \structure{structure}. Und jetzt nochmal ähnlich in Blockform:
	
	\begin{block}{Titel}
		Text kommt hier rein
	\end{block}
	
	\begin{alertblock}{Titel}
		Text kommt hier rein
	\end{alertblock}
	
	\begin{exampleblock}{Titel}
		Text kommt hier rein
	\end{exampleblock}
\end{frame}

\section{Formeln und Formeln im Block}
\begin{frame}{Das Rasch-Modell}
	\begin{equation}
		\frac{exp(\theta-\beta)}{1+exp(\theta-\beta)}
	\end{equation}
\end{frame}

\begin{frame}{Das Rasch-Modell}
	
	\begin{block}{Die Formel}
		\LARGE
			\begin{equation}
			\frac{exp(\theta-\beta)}{1+exp(\theta-\beta)}
		\end{equation}
	\end{block}
\end{frame}

\section{Gedichte}
\begin{frame}{Die Geschichte vom Suppen-Kaspar}
	 \begin{verse}
	 	\footnotesize
	 	\begin{columns}[T]
	 		\begin{column}{0.49\textwidth}
 				Der Kaspar, der war kerngesund,\\
	 			Ein dicker Bub und kugelrund,\\
	 			Er hatte Backen rot und frisch;\\
	 			Die Suppe aß er hübsch bei Tisch.\\
	 			Doch einmal fing er an zu schrei’n:\\
	 			„Ich esse keine Suppe! Nein!\\
	 			Ich esse meine Suppe nicht!\\
	 			Nein, meine Suppe ess’ ich nicht!“\linebreak
	 				 			
	 			Am nächsten Tag, — ja sieh nur her!\\
	 			Da war er schon viel magerer.\\
	 			Da fing er wieder an zu schrei’n:\\
	 			„Ich esse keine Suppe! Nein!\\
	 			Ich esse meine Suppe nicht!\\
	 			Nein, meine Suppe ess’ ich nicht!“
	 		\end{column}
	 		\begin{column}{0.49\textwidth}
			 	Am dritten Tag, o weh und ach!\\
	 			Wie ist der Kaspar dünn und schwach!\\
	 			Doch als die Suppe kam herein,\\
	 			Gleich fing er wieder an zu schrei’n:\\
	 			„Ich esse keine Suppe! Nein!\\
	 			Ich esse meine Suppe nicht!\\
	 			Nein, meine Suppe ess’ ich nicht!“\linebreak
	 			
	 			Am vierten Tage endlich gar\\
	 			Der Kaspar wie ein Fädchen war.\\
	 			Er wog vielleicht ein halbes Lot —\\
	 			Und war am fünften Tage tot.\\
	 		\end{column}
	 	\end{columns}	 	
	 \end{verse}
\end{frame}

\section{Eigene Quote-Umgebung}
\begin{frame}
	\begin{authquote}[Steven Erikson - Memories of Ice]
		Outside the city's west wall, close to the shoreline's broken, jagged edge, a lazy swirl of dust rose from the ground, took form. Tool slowly settled the flint sword into its shoulder hook, his depthless gaze ignoring the abandoned shacks to either side and fixing on the massive stone barrier before him.
		
		Dust on the wind could rise and sweep high over this wall. Dust could run in the streams through the rubble fill beneath the foundation stones. The T'lan Imass could make his arrival unknown.
		
		But the Pannion Seer had taken Aral Fayle. Toc the Younger. A mortal man... who had called Tool friend. He strode forward, hide-wrapped feet kicking through scattered bones.
		
		The time had come for the First Sword of the T'lan Imass to announce himself.
	\end{authquote}
\end{frame}

\begin{frame}[plain]{Folie plain}
	Hier ist nichts.
\end{frame}

\section{Abschnittsfolie}

\begin{frame}
	\sectionpage
\end{frame}

\subsection{Abschnittsfolie mit Unterabschnitt}
\begin{frame}
	\sectionpage
\end{frame}

\subsection{Unterabschnittsfolie}
\begin{frame}
	\subsectionpage
\end{frame}
	
	
	
	
\end{document}
